(1) 量子比特等先验概率制备为 $|\psi_0\rangle=|0\rangle$ 或 $|\psi_1\rangle=\cos\theta|0\rangle+\sin\theta|1\rangle$，其中 $0<\theta<\pi/2$。在基 $|0\rangle,|1\rangle$ 中写出密度矩阵 $\rho_0,\rho_1$ 以及差算符 $\Delta=\rho_0-\rho_1$。
(2) 求 $\Delta$ 的两个本征值、迹范数 $\|\Delta\|_1$，根据Helstrom定理计算最优成功率和最小错误概率。
(3) 描述实现Helstrom最优判别的投影测量，并讨论 $\theta\to0^+$ 与 $\theta\to\pi/2^-$ 的极限。
(4) 若改要求绝对不会给出错误结论而允许“不确定”第三结果，说明等先验两纯态无歧义判别可达的最小不确定概率，并与第(2)问的任务区别。
